% Generated draft: requires mathematical and source review.
\documentclass[11pt]{article}
\usepackage[margin=25mm]{geometry}
\usepackage{fontspec}
\setmainfont{Latin Modern Roman}
\usepackage{amsmath,amssymb,mathtools,unicode-math}
\setmathfont{Latin Modern Math}
\usepackage{enumitem,needspace,xurl,hyperref}
\setlength{\parindent}{0pt}
\setlength{\parskip}{5pt}
\setlength{\emergencystretch}{3em}
\newcommand{\N}{\mathbb{N}}
\newcommand{\Z}{\mathbb{Z}}
\newcommand{\Q}{\mathbb{Q}}
\newcommand{\R}{\mathbb{R}}
\newcommand{\C}{\mathbb{C}}
\newcommand{\F}{\mathbb{F}}
\newcommand{\A}{\mathbb{A}}
\newcommand{\PP}{\mathbb{P}}
\newcommand{\E}{\mathbb{E}}
\newcommand{\eps}{\varepsilon}
\newcommand{\dd}{\,\mathrm{d}}
\begin{document}
\section*{UnsolvedMath ID: 20000754}
\textbf{Source problem code:} catalogue entry
\subsection{Definitions and mathematical statement}
Let \(S=\{s_1<s_2<\cdots\}\) be a specified increasing sequence of distinct positive integers. A bounded-diameter consecutive triple means indices \(j\) with \(s_{j+2}-s_j\leq C\), for a constant independent of \(j\). C. Elsholtz asks for infinitely many such triples in sequences resembling the primes but more manageable, especially
\[
S=\{a^2+b^2>0:a,b\in\Z\}.
\]
K. B. Ford notes the elementary clusters \(n^2,n^2+1,n^2+4\). Even when these three displayed values are not themselves consecutive in \(S\), their common interval contains three members and hence a consecutive triple of span at most four. J. B. Friedlander suggests excluding such degenerate representations by requiring the two squares to have comparable size; S. W. Graham and T. D. Wooley instead suggest obtaining ``many'' gaps. No comparability constant or numerical abundance threshold is specified. R. C. Vaughan proposes the broader setting of norm-form values, without choosing a field or form.

A precise comparability convention is to require a representation \(m=a^2+b^2\) with \(a,b>0\) and \(1/2\leq a/b\leq2\). Then the squares themselves differ by at most a factor of four. A primitive representation additionally requires \((a,b)=1\); that is an extra strength of the attached aid, not an original requirement.

The aid reports a fixed arithmetic progression of positive integers \(n\), defined by Chinese-remainder congruences, such that for all sufficiently large members of that progression the three numbers
\[
5n^2+5,\qquad5n^2+20,\qquad5n^2+45
\]
are exact consecutive terms of the full sequence \(S\). Their representations are
\[
5n^2+5k^2=(n+2k)^2+(2n-k)^2,\qquad k=1,2,3,
\]
and the aid claims each is primitive and \(2\)-comparable in the coordinate-ratio convention. Thus the two gaps are exactly \(15\) and \(25\), with span \(40\). It reports at least \(c\sqrt X\) such triples with largest member at most \(X\), for some fixed \(c>0\) and all sufficiently large \(X\). This makes one possible meaning of ``many'' explicit, without asserting a positive proportion of all sums of two squares or an optimal count. The actual modulus and residues of the progression are not supplied, so this record does not reproduce the full construction.

The aid also reports an orthogonality-based extension: for every fixed positive-definite integral binary quadratic form \(Q\), the increasing sequence of its distinct positive represented values has infinitely many consecutive triples of bounded diameter depending on \(Q\). The summary does not assert the same primitive/comparable representation properties for every such form. It explicitly does not cover general higher-degree norm forms. These are supplied unreviewed partial claims, not a resolution of every prime-like sequence or every interpretation of abundance.
\subsection{Short English statement}
Seek infinitely many bounded-span consecutive triples in prime-like sequences, with nondegenerate representations or quantitative abundance as refinements. The supplied aid reports a primitive comparable sum-of-two-squares construction with exact gaps \(15,25\), plus a binary-quadratic-form extension, without treating all norm forms.
\subsection{Sources}
\begin{itemize}
\item C. Elsholtz and the comments of K. B. Ford, J. B. Friedlander, S. W. Graham, T. D. Wooley, and R. C. Vaughan, Problem 19, p. 5, Angel Kumchev's notes from the December 2005 ARCC workshop \emph{Gaps Between Primes}: \url{https://aimath.org/WWN/primegaps/WorkshopProblems.pdf}.
\item Local registry \texttt{data/unsolvedmath/problems.json}, permanent ID \(20000754\), exact source code \texttt{catalogue entry}; canonical \texttt{aim-analytic-number-theory-notes.json}, notes[117].
\item The attached research attempt supplies the CRT construction and binary-form extension without a complete external citation or the progression data. It is explicitly machine-generated, not peer reviewed, and requires independent verification before citation.
\end{itemize}
\end{document}
